\documentclass{article}
\usepackage{tikz}
\usetikzlibrary{angles, quotes}
\usepackage{amssymb}
\usepackage{pgfplots}
\pgfplotsset{width=10cm,compat=1.9}
\usepgfplotslibrary{statistics}
\usepackage[left=2.54cm,right=2.54cm,top=2.54cm,bottom=2.54cm]{geometry}
\usepackage{tasks} %Paquete necesario para la producción de items horizontales para algunos ejercicios de tarea empleando el comando /begin{multicols}}{#}.../end{multicols} para ayudarnos a reducir las páginas
\usepackage{pdflscape} %comando que permite cambiar a orientación horizontal las hojas%
\usepackage{soul} %Paquete para permitir la compilación de acentos usando el comando {\hl{}}}\\
\sethlcolor{yellow(munsell)} %comando para definir el color para el subrayado usando el comando \hl
\usepackage[utf8]{inputenc}
\usepackage[spanish]{babel}
\usepackage{icomma}
\usepackage{ragged2e}
\usepackage{siunitx}
\usepackage{url}
\usepackage[colorlinks=true, urlcolor=blue,  linkcolor=black, citecolor=green]{hyperref}
\usepackage{pdfpages}
\usepackage{blindtext} %Paquete que genera texto ficticio (dummy text) con el comando: \blindtext
\usepackage{cancel} %in the preamble gives you four different modes of striking through: \cancel{text to cancel} draws a diagonal line (slash) through its argument, \bcancel{text to cancel} uses the negative slope (a backslash), \xcancel{text to cancel} draws an X (actually \cancel plus \bcancel), \cancelto{〈value〉}{〈expression〉} draws a diagonal arrow through the 〈expression〉pointing to the 〈value〉 (math-mode only)
\usepackage{csquotes}
\usepackage{afterpage}
\usepackage{parskip} 
\usepackage{float}
\usepackage{enumitem}
\usepackage{multicol}%Paquete que permite la creación aislada de columnas de texto. De esta forma se reduce la cantidad de páginas en nuestro documento
\usepackage{lipsum} %Paquete que genera texto de relleno al igual que \usepackage{blindtext}, se usa con el comando \lipsum[#-#] los #(hashtags) delimitan el número de páginas que deseamos ocupar del paquete lipsum para usarlos como relleno. 

% Definir el comando personalizado para la nota
\newcommand{\mynote}[1]{%
\begin{tikzpicture}[baseline=-0.75ex]
\node[draw, circle, fill=Apple Green!60, inner sep=2pt] (note) {\textbf{Nota:}};
\end{tikzpicture}%
\ \textit{#1}%
}

\newenvironment{Figura}
{\par\medskip\noindent\minipage{\linewidth}}
{\endminipage\par\medskip}
\usepackage{caption}
\usepackage[
backend=biber,
sorting=none,
url=true
]{biblatex} %bibliografía
\addbibresource{biblio.bib}
% Establecer el espacio entre las entradas de la bibliografía
\setlength{\bibitemsep}{1\baselineskip} % Puedes ajustar el valor según tus preferencias
\usepackage{amsmath,amsthm,amssymb,amsfonts}
\usepackage{pifont} %Permite la compilación del comando \xmark: es el símbolo de la tachita
\usepackage{mathtools} %permite la compilación de símbolos para matrices
\usepackage{empheq} %Paquete que se relaciona con \usepackage[most]{tcolorbox} para la creación de cuadros/cajas de colores para delimitar los resultados matemáticos o líneas de texto usando la línea de comando: \begin{empheq}[box={\mymath[colback=orange(webcolor)!70,drop lifted shadow]}]{equation*} \end{empheq}
\usepackage[most]{tcolorbox} %Paquete que se relaciona con \usepackage{empheq} para la creación de cuadros/cajas de colores para delimitar los resultados matemáticos o líneas de texto
\renewcommand{\qedsymbol}{$\blacksquare$}
\usetikzlibrary{positioning,decorations.pathreplacing} %Paquete que permite la compilación de llaves para cuadros sinópticos (brace diagram) 
\usepackage{schemata} %The schemata package is designed just for these brace diagrams
\usepackage{booktabs}

\newcommand\AB[2]{\schema{\schemabox{#1}}{\schemabox{#2}}} %comando o paquete necesario para crear cuadros sinópticos

\newtcbox{\mymath}[1][]{%
nobeforeafter, math upper, tcbox raise base,
enhanced, colframe=blue!30!black,
colback=blue!30, boxrule=1pt,
#1} %Comando importante para encasillar los resultados matemáticos en cajas de diferentes colores, se relaciona con los paquetes \usepackage{empheq} y \usepackage[most]{tcolorbox} 

\newenvironment{sysmatrix}[1]
{\left(\begin{array}{@{}#1@{}}}
{\end{array}\right)}
\newcommand{\ro}[1]{%
\xrightarrow{\mathmakebox[\rowidth]{#1}}%
}
\newlength{\rowidth}% row operation width
\AtBeginDocument{\setlength{\rowidth}{3em}}  %Comando importante para laproducción de lineas de operaciones entre matrices, método de Gauss o Gauss Jordan se relaciona con \newenvironment{sysmatrix}

%formato para cambiar el horario a español
\usepackage[spanish]{datetime2}
\DTMsetdatestyle{spanish}

\renewcommand{\today}{\DTMdisplaydate{\the\year}{\the\month}{\the\day }{-1}}
%%%%%%%%%%%%%%%%%%%%%%%%%%%%%%%%%%%%%%%%%

\usepackage{datetime}
\newcommand{\mycurrenttime}{\xxivtime}

%%%%%%%%%%%%%%%%%%%%%%%%%%%%%%%%%%%%%%%%%%%%%%%%%%%%%%%%%%%%%%%%%%%%%%%%%%%%%%%%%%%%%%%%%%%%%%%%%%%%%%%%%%%%%%%%%%%%%%%%%%%%%%%%%%%%%%%%%%%%%%%%%


%%%%%%%%%%%%%Caja de color para el título%%%%%%%%%%%%%%%%%%%%%%%%%%%%%%%%%%%%%%%%%%%%%%%%%%%%%%

\definecolor{myframecolor}{RGB}{1, 45, 75} % Prussian Blue
\definecolor{myboxcolor}{RGB}{0, 107, 92} % Apple Green

\newtcolorbox{mybox}{
enhanced,
colback=myboxcolor!7, % Color de fondo del cuadro
colframe=myframecolor, % Color del marco del cuadro
arc=0pt,
boxrule=1pt,
borderline west={2mm}{-10mm}{myframecolor}, % Borde en el lado izquierdo
sharp corners=southwest,
width=\linewidth,
before=\par\vspace{\bigskipamount}, % Espacio antes del cuadro
after=\par\vspace{\bigskipamount} % Espacio después del cuadro
}


\newcommand{\euler}{e} %Comando para producir letra e de euler
\newenvironment{remark}{\par\vfill\footnotesize % Vertical white space above the remark and smaller font size
\begin{list}{}{
	\leftmargin=80pt % Indentation on the left
	\rightmargin=60pt}\item\ignorespaces % Indentation on the right
\makebox[-2.5pt]{\begin{tikzpicture}[overlay]
		\node[draw=Horizon!60,line width=2.5pt,circle,fill=Horizon!25,font=\sffamily\bfseries,inner sep=4pt,outer sep=0pt] at (-19pt,5pt){\textcolor{Horizon}{Nota}};\end{tikzpicture}} % Blue Nota in a circle
\advance\baselineskip -1pt}{\end{list}\vskip5pt} % Tighter line spacing and white space after remark
\usepackage{graphicx}

\usepackage{titling}

\input{colores.tex}

%%Producir signos de cita textual``Asignatura''

\newcommand{\xmark}{\ding{55}} %%%comando que genera la tachita

\newenvironment{MyColorPar}[1]{%
\leavevmode\color{#1}\ignorespaces%
}{%
}%

%%%%%%%%%%%%%%%%%%% Título de la tarea, Nombre de alumno

\title{Clase \# 32 - Estado Sólido}
\author{\emph{JA} $\boldsymbol{\mid}$}
\date{\today}

%%%%%%%%%%%%%%%%%%% Datos de la Materia

\usepackage{fancyhdr}
\fancypagestyle{plain}{%  the preset of fancyhdr 
\fancyhf{} % clear all header and footer fields
\fancyfoot[R]{
\includegraphics[width=2cm]{LogoFCFMBUAP (1).png}}
\fancyfoot[L]{{\bfseries{\thedate{}}} a las {\bfseries{\mycurrenttime{}}} horas{} (GMT-6, H. Puebla de Zaragoza, Pue)}
\fancyhead[L]{Estado sólido}
\fancyhead[R]{\theauthor}
}
\makeatletter
\def\@maketitle{%
\newpage
\null
\vskip 1em%
\begin{center}%
\let \footnote \thanks
{\LARGE \@title \par}%
\vskip 1em%
%{\large \@date}%
\end{center}%
\par
\vskip 1em}
\makeatother

\usepackage{lipsum}  
\usepackage{cmbright}


%%%%%%%%%%%%%%%%%%%%%%%%%%%%%%%%%%%%%%%%%%%%%%%%%%%%%%%%%%%%%%%%%%%%%%%%%%%%%%%%%%%%%%%%%%%%%%%%%%%%%%%%%%%%%%%%%%%%%%%%%%%%%%%%%%%%%%%%%%%%%%%%%%%%%%%%%%%%%%%%%%%%%%%%%%%%%%%%%%%%%%%%%%%%%%%%%%%%%%

\begin{document}



\maketitle


%%%%%%%%%%%%%%%%%%% Datos del profesor, campus y aula

\begin{mybox}
\noindent\begin{tabular}{@{}ll}
	{\bfseries{Profesor}} & Dra. Ana Lilia González Ronquillo\\
	\textcolor{prussianblue}{\bfseries{Campus}} / \textcolor{trueblue}{\bfseries{Facultad}}  & \textcolor{prussianblue}{\bfseries{C.U. BUAP}} /  \textcolor{trueblue}{\bfseries{FCFM}} \\
	\textcolor{prussianblue}{\bfseries{Edificio}} /  \textcolor{trueblue}{\bfseries{Salón}}    & 	\textcolor{prussianblue}{\bfseries{FM9}} /  \textcolor{trueblue}{\bfseries{301}}   
\end{tabular} 
\end{mybox} \vspace{1cm}

\tableofcontents

\vspace{2cm}

Tenemos:

\begin{equation}
\begin{split}
 E  &= \int\limits_{0}^{\omega D} \frac{  \hbar \omega g(\omega) }{ e^{\hbar \omega / K_{B}T } -1  dw} = \frac{L}{\pi \nu} \int\limits_{0}^{\omega{D}} \frac{ \hbar \omega }{ e^{\hbar \omega / K_{B}T} -1  } d\omega
\end{split}
\end{equation}


EN 3D cada valor de $k$ ocupa un volumen:


\[ v_{k} = \left( \frac{2\pi}{L} \right)^{3} \]


¿Cuál es el nimero de modos entre el intervalo? 

Es una muestra cúbica, un pedazo de material cúbico, sin embargo la estructura de adentro no la tomamos en cuenta:


Entonces:


\[  g(k) = \frac{ \text{No. de modos en un intervalo}  }{  \text{ volumen del intervalo } } = \frac{ 1 }{ \left(  \frac{2\pi}{L} \right)^{3} } = \left( \frac{L}{2\pi} \right)^{3} = \frac{V}{8\pi^{3}}\]

\begin{equation}
\begin{split}
 g(k) \to g(\omega) \to E &\to C_{v}\\
  k &= \sqrt{ k_{x}^{2} + k_{y}^{2} + k_{z}^{2} }
\end{split}
\end{equation}

A manera de aproximación se usa el nímero total $N$ de valores de $K$ dentro de una esfera de radio $\mathbf{\mathcal{K}}$:

\[N= \frac{ V_{esf} }{V_{K}} = \frac{V}{ 8\pi^{3}} \left(  \frac{4\pi k^{3}}{3} \right)   \]

\begin{figure}[H]
    \centering
	    \includegraphics[width=0.8\textwidth]{1.png} % ajusta el ancho según convenga
    \caption{Descripción clara y concisa de la figura.}
    \label{fig:mi_figura} % opcional: para usar con \ref{fig:mi_figura}
\end{figure}

\[  k = \left| \mathbf{\mathcal{K}} \right| = \sqrt{  k_{x}^{2} + k_{y}^{2} + k_{z}^{2}  } \]

La densidad $g(\omega)$ es:


\begin{figure}[H]
    \centering
	    \includegraphics[width=0.8\textwidth]{2.png} % ajusta el ancho según convenga
    \caption{Descripción clara y concisa de la figura.}
    \label{fig:mi_figura} % opcional: para usar con \ref{fig:mi_figura}
\end{figure}

Considerando una expresión similar de densidad de cada modo:

\[  g(\omega) = \frac{V }{2\pi^{2}} \left\{  \left(  \frac{ \omega }{ v_{L} }  \right)^{2} \left( \frac{ 1 }{v_{L}}  \right) + \left(\frac{\omega}{v_{T1}} \right)^{2} \left(  \frac{1}{v_{T1}} \right) + \left( \frac{\omega}{v_{T2}} \right)^{2} \left( \frac{1}{v_{T2}} \right) \right\}  \]

¿Qué pasaría en el caso de BCC y FCC, los modos longitudinales y transversales pueden tener diferentes velocidades? Se pueden propagar de igual forma los longitudinales que los tran{}versales.

La estructura determina si los modos se propagan más rápido o más lento. En la CS todo es exactamente igual. Ambos modos tienen la misma velocidad. Tiene que ver con la densidad de átomos. SIiexiste más materia hay más probabilidad de que se propague la onda. 

Velocidad de grupo: velocidad con la que propaga energía. 


En su modelo, Debye asumió por simplicidad: $v_{L} = v_{T{1}} = v_{T2}=v_{0}$, velocidad de la onda de sonido en el material.

Por lo que:

\[ \textcolor{Cerulean}{\boxed{\textcolor{onyx}{g(\omega) = \frac{ 3V\omega^{2} }{ 2\pi^{2}v_{0}^{3} }} }}   \]

Sustituyendo en la integral:


\[  3N = \int\limits_{0}^{\omega{D}} g (\omega) d\omega\]

Sustituyendo en la integral, tenemos: 

\[  \textcolor{orange(colorwheel)}{\boxed{\textcolor{onyx}{\omega_{D} = \left( \frac{6\eta \pi^{2}}{V}  \right)^{\frac{1}{3}} v_{0}} }}   \]

$\eta =$ densidad

\begin{figure}[H]
    \centering
	    \includegraphics[width=0.8\textwidth]{3.png} % ajusta el ancho según convenga
    \caption{Descripción clara y concisa de la figura.}
    \label{fig:mi_figura} % opcional: para usar con \ref{fig:mi_figura}
\end{figure}


Tenemos en:


\begin{itemize}
\item[$\bullet$] En 1D: $g(\omega) = \frac{L}{\pi v_{0}} = cte = cte \omega^{1-1} = cte \omega^{0}$ 
\item[$\bullet$] En 2D: $g(\omega) \propto \omega^{D-1}= \omega^{2-1} = \omega$ 
\item[$\bullet$] En 3D: $g(\omega) \propto \omega^{D-1}\omega= \left( 6 \eta \pi^{2} \right)^{\frac{1}{3}} v_{0}$ 
\end{itemize}

\begin{figure}[H]
    \centering
	    \includegraphics[width=0.8\textwidth]{4.png} % ajusta el ancho según convenga
    \caption{Descripción clara y concisa de la figura.}
    \label{fig:mi_figura} % opcional: para usar con \ref{fig:mi_figura}
\end{figure}


\begin{figure}[H]
    \centering
	    \includegraphics[width=0.5\textwidth]{5.png} % ajusta el ancho según convenga
    \caption{Descripción clara y concisa de la figura.}
    \label{fig:mi_figura} % opcional: para usar con \ref{fig:mi_figura}
\end{figure}

Para bajas $T^{\prime}s$:

\[ T<< \theta_{D}/10  \] el error es muy pequeño si se considera: $\int\limits_{0}^{\infty}$


\[ E \approx \frac{ 9NK_{B}T^{4} }{ \theta_{D}^{3} } \underbrace{\int\limits _{0}^{\infty} \frac{x^{3}}{e^{x}-1} dx}_{\boxed{\frac{\pi^{4}}{15}}, \text{funciones especiales}}	= \frac{ 3\pi^{4}NK_{B} }{5} \left(  \frac{ T^{4} }{\theta_{D}^{3}} \right)\]

\begin{figure}[H]
    \centering
	    \includegraphics[width=0.5\textwidth]{6.png} % ajusta el ancho según convenga
    \caption{Descripción clara y concisa de la figura.}
    \label{fig:mi_figura} % opcional: para usar con \ref{fig:mi_figura}
\end{figure}

\begin{figure}[H]
    \centering
	    \includegraphics[width=0.5\textwidth]{7.png} % ajusta el ancho según convenga
    \caption{Descripción clara y concisa de la figura.}
    \label{fig:mi_figura} % opcional: para usar con \ref{fig:mi_figura}
\end{figure}



A altas temperaturas los 3 modelos tienden a lo mismo.  A muy bajas temperaturas \textbf{debye} es bastante aproximado. Cuando no es tan alta ni tan baja, tomar aleatoriedad.

\begin{figure}[H]
    \centering
	    \includegraphics[width=0.8\textwidth]{8.png} % ajusta el ancho según convenga
    \caption{Descripción clara y concisa de la figura.}
    \label{fig:mi_figura} % opcional: para usar con \ref{fig:mi_figura}
\end{figure}


Propiedades térmicas de las estructuras cristalinas de los átomos

$KCl$ es un aislante no hay electrones de conducción.  $Cu$ es un conductor, el modelo de \textbf{Debye} no contempla el calor electrónico específico (\underline{electronic specific heat}.)

La siguiente clase trabajaremos con electrones. Es un material que tiene una naturaleza diferente.

\[ C_{v} \quad \to \quad \text{oscilaciones de la red (fonones)} \]

\end{document}

